\section{Conclusion}
\label{sec:conclusion}

%% 目的：兩個結論各一句。138 那個數字原本在 §5.5 是帶著punchline的（「而實際上只有 14 個這種錯誤」），
%%       結論卻只留前半——讀者會以為修 138 個 promise FN 是另一條可走的路。
In a gated composite metric the published field weights do not say what a marginal improvement
contributes, and two consequences of the scorer, under schema-respecting predictions, point to
opposite allocations of effort. A scored
class carries $w_f/|C_f|$ whatever its support, so \textsc{Misleading} is worth as much as $138$
promise false-negative repairs would return at our prices --- against the 14 such errors that exist. And because a gate
error is charged in the fields it gates --- twice over in evidence, where N/A is scored --- the
nominally $0.20$-weighted gate has an effective weight of $0.5216$ in that direction ($0.4268$ in
the other). The first needs only a
task overview, the second also an operating point: both precede a first submission.

%% 目的：裁決。原本這一段宣告「We settled the conflict by measurement」之後就再也沒說判決是什麼，
%%       後面全是證據，最後收在釋出程式碼——只讀結論的人會拿不到 §1 承諾要交付的那個東西。
We settled the conflict by measurement, and the verdict is that a price is not an invoice. The
gate's gains were reachable by ordinary subset selection; on the rare class the only unaided
precision we can measure is $0/35$, and the one confirmed test placement came out of staking five
items against leaderboard feedback --- a search, not a detector, and not something we could plan
effort around. Only the upstream work showed a gain on validation, and with supports of 1 and 3
the rare class yields no general rule.

%% 目的：證據與誠實的負結果。原句「reproduces all four field macro-F1 scores exactly」和 §4.2 現在說的
%%       「重現接近恆等式」衝突；「our detector flagged one and missed two」則把 §7.1 和 §8.2 兩度拒絕
%%       這樣讀的數字重新讀成偵測器成績。
The cascade model has no tuned parameters: every rate is counted on the $1{,}000$ validation
items, so reproducing the four field macros is arithmetic, not evidence. The demand we put on it is
transfer, checked on leaderboard-informed predictions: its counted real-class gate-open rates carry to the released test labels within $0.05$, though its
predicted quality macro does not (Section~\ref{sec:disc:limits}). We reached the rare
class only in part: on validation the repick puts 35 items in the class and none is right, and on
test we staked five and the labels confirm one, a placement worth $+0.0292$ of the total.

%% 目的：估計面，並收在 §1 承諾的三條建議上（原本收在釋出程式碼，URL 摘要和 §4.4 已經各給過一次）。
The estimation side is bounded twice: support, not split size, is what a macro-F1 term resolves,
and a half-width four times the seed spread is why one of our own deployed decisions sits below
what validation can settle. We close with three
additions to what a task overview reports --- a conditional-on-gate downstream score,
per-partition rare-class supports, and room in the schema for greenwashing that accompanies no
promise --- so the next participant can read the weights and know what they are being paid for.
